%% kapa: exam=1 round=20 year=2009 period=1 form=A subject=civil_law subject_ko=민법 seq=32 number=32
%% source: Q-Net 감정평가사 기출문제 (cst00309) artlSeq=5024178 file=extracted/1008/5024178/2009년 제20회 감정평가사 1차시험 기출문제/1교시.hwp
%% conversion: kordoc 4.24.51 -> tools/gen1.py (auto); review=proofread-sonnet; flags=none
\begin{question}{32}
甲과 乙은 乙소유의 건물에 대하여 전세금 1억원에 전세권을 설정하기로 하고 등기도 경료하였다. 다음 중 옳지 \uline{않은} 것은? (다툼이 있는 경우 판례에 의함)
\begin{choices}
  \item 甲은 전세권 존속 중이라도 장래에 그 전세권이 소멸하는 경우에 전세금반환채권이 발생하는 것을 조건으로 그 장래의 조건부 채권을 유효하게 양도할 수는 있다.
  \item 전세권이 소멸한 경우, 乙의 전세금반환의무와 甲의 목적물 인도 및 전세권말소등기에 필요한 서류의 교부의무는 동시이행의 관계에 있다.
  \item 甲이 乙에 대한 기존의 채권 1억원으로 전세금지급에 갈음한 경우에는 전세권은 성립되지 아니한다.
  \item 만일 甲과 乙이 전세금반환채권을 담보할 목적으로 甲ㆍ乙ㆍ丙 3자간의 합의에 의해 丙명의로 전세권설정등기를 경료하였더라도 이 역시 유효하다.
  \item 전세권이 성립한 후 목적물의 소유권이 乙에서 丁으로 이전된 경우 乙은 원칙적으로 전세권설정자의 지위를 상실하여 전세금반환의무를 면한다.
\end{choices}
\end{question}
\answer{3}
